%% Source-based comparison; no new transport calculations.
%% External MOCASSIN sources are preserved for comparison.
\documentclass[english,a4paper]{my_memo}
\synctex=1
\usepackage{amsmath}
\usepackage{amssymb}
\usepackage{microtype}
\usepackage{booktabs}
\usepackage{longtable}
\usepackage{array}
\usepackage{enumitem}
\usepackage{placeins}
\hypersetup{hidelinks}
\newcommand{\code}[1]{\texttt{#1}}
\newcommand{\MoCHII}{\textsc{MoCHII}}
\newcommand{\Mocassin}{\textsc{Mocassin}}
\newcommand{\HI}{H\,\textsc{i}}
\newcommand{\HII}{H\,\textsc{ii}}
\newcommand{\Te}{T_{\rm e}}
\newcommand{\nelec}{n_{\rm e}}
\newcolumntype{P}[1]{>{\raggedright\arraybackslash}p{#1}}
\newcommand{\Source}{\noindent\textbf{Source reading.}\ }
\newcommand{\Measured}{\noindent\textbf{Quoted measurement.}\ }
\newcommand{\Inferred}{\noindent\textbf{Inference.}\ }
\newcommand{\mc}[2]{\textsf{M:}\,\allowbreak\texttt{\detokenize{#1}}:\allowbreak#2}
\newcommand{\ms}[2]{\textsf{S:}\,\allowbreak\texttt{\detokenize{#1}}:\allowbreak#2}
\newcommand{\old}[2]{\textsf{O:}\,\allowbreak\texttt{\detokenize{#1}}:\allowbreak#2}
\newcommand{\both}[1]{\multicolumn{2}{P{\dimexpr.52\linewidth+2\tabcolsep\relax}}{#1}}
\begin{document}
\title{MOCASSIN and MoCHII:\\Three Monte Carlo Radiation-Transfer Algorithms}
\author{Kwang-Il Seon}
\date{\docmoddate}
\rhead{\bfseries MOCASSIN and MoCHII: Monte Carlo algorithms}
\maketitle
\begin{abstract}
This memo compares the connected Monte Carlo calculation in the local
MOCASSIN 2.02.73.2 source tree with two MoCHII algorithms: the v1.00
leaf-emissivity method, represented by the same option in v1.10, and
v1.10 absorption re-emission. All three iterate transport and local gas
balance. Their principal differences concern packet weights, spectral
sampling, radiation estimators, and how diffuse emission enters packet
histories. Absorption re-emission moves MoCHII closer to MOCASSIN in the
location and timing of emission, but retains MoCHII's continuous-energy
estimators, weighted histories, and explicit emission remainder. It is
not MOCASSIN's indivisible-packet algorithm. Source readings, previously
measured MoCHII results, and mathematical inferences are labeled
separately. No MOCASSIN calculation or new transport calculation was run
for this memo. Atomic-data comparisons are excluded.
\end{abstract}
\tableofcontents

\section{Scope, evidence, and summary}
\label{sec:scope}
\Source The three columns below mean:
\begin{description}[leftmargin=2.1em,style=nextline,itemsep=2pt]
\item[A: MOCASSIN] The inspected tree
\path{~/MOCASSIN/mocassin-mocassin.2.02.73.2/source/}, identified by the author on November 14, 2026 as the original-author code, with the five supplied \path{~/MOCASSIN/error_fix/} files already incorporated.
\item[B: v1.00 algorithm] The default leaf-emissivity algorithm, evaluated
in the current v1.10 source with
\code{par\%diffuse\_launch}\allowbreak\code{='leaf\_emissivity'}.
\item[C: v1.10 new option] The same code with
\code{par\%diffuse\_launch}\allowbreak\code{='absorption\_reemission'} and transported
recombination radiation, using the current capped r7 rule.
\end{description}
References \textsf{M:}, \textsf{S:}, and \textsf{O:} identify files in that
MOCASSIN directory, \path{MoCHII_v1.10/src/}, and
\path{MoCHII_v1.00/src/}, respectively. MOCASSIN line numbers refer to the original-author tree rechecked on
November 14, 2026; the MoCHII citations retain their earlier source basis. An isolated routine is not counted as an active
algorithm without following its caller.

\Source The five files in \path{~/MOCASSIN/error_fix/} match the
corresponding files in source A. ``Original-author code'' identifies this
specified tree, not an independently authenticated pristine release;
its supplied files also contain dated corrections. Its driver calls
input, grid/radiation setup, and \code{MCIterationDriver}, which enters
\code{iterateMC} (\mc{mocassin.f90}{65--83,\allowbreak 124--126};
\mc{../Makefile}{41--45}; \mc{iteration_mod.f90}{12--34}).
The separate \path{~/RT_Codes/MOCASSIN/mocassin-mocassin.2.02.73.2r/source/}
tree is the user's unfinished error-fix and optimization work, including
fixed-seed input, explicit-face traversal, and revised MPI/convergence
code; it is not source A. Differences relevant here are stated where
they affect the description. No equivalence of the two executables is
assumed.

\Source B denotes algorithm equivalence, not universal executable or byte
identity. Direct source comparison finds identical common spectral,
QMC, scheduling, ion-balance, thermal, and stopping modules. The diffuse
launcher uses \code{dif\_launch=>dif\_ch} in B and retains the original
leaf/channel sampling (\ms{diffuse_mod.f90}{250--252,\allowbreak 399--406,\allowbreak 677--763};
\old{diffuse_mod.f90}{400--486}). The current direct-flight implementations
also agree, including stable thin-optical-depth arithmetic
(\ms{raytrace_amr.f90}{680--755,\allowbreak 803--869};
\old{raytrace_amr.f90}{491--566,\allowbreak 614--680}). This is a comparison of methods,
not an archived v1.00 executable.

\noindent\textbf{Evidence labels.} \emph{Source reading} describes
inspected executable expressions and connected calls.
\emph{Quoted measurement} reproduces an identified existing report
without re-running its calculation. \emph{Inference} is a consequence
of the algorithms, not a measured cross-code result. Detailed sections
supply file/line evidence for every aspect in the summary. Merged B/C
cells mean that the stated operation is shared.

\clearpage
\begingroup\small
\setlength{\tabcolsep}{4pt}
\renewcommand{\arraystretch}{1.12}
\begin{longtable}{@{}P{.15\linewidth}P{.27\linewidth}P{.26\linewidth}P{.26\linewidth}@{}}
\caption{Algorithm summary. A is the inspected local MOCASSIN tree;
B and C are the two MoCHII choices. Source locations are in
Sections~\ref{sec:same}--\ref{sec:bc}.}\label{tab:summary}\\
\toprule Aspect & MOCASSIN (A) & v1.00 algorithm (B) & v1.10 re-emission (C)\\
\midrule\endfirsthead
\toprule Aspect & MOCASSIN (A) & v1.00 algorithm (B) & v1.10 re-emission (C)\\
\midrule\endhead
\bottomrule\endfoot
Packet & Fixed normalization within an ordinary history; no evolving weight &
\both{Luminosity normalization times a dimensionless weight; physical photon energy stored separately}\\
Spectrum & Stored discrete frequency from a CDF &
\both{Continuous energy from spectral/channel samplers; line components can have prescribed energies}\\
Sources & Explicit source loops and integer allocations &
\both{Component drawn in proportion to band luminosity; optional separate FUV sampling}\\
Directions & Intrinsic pseudorandom generator; geometry dependent &
\both{Pseudorandom or scrambled Sobol launch coordinates; geometry dependent}\\
Direct flight & Random exponential optical depth and crossed lengths &
\both{Analytically attenuated chosen launch ray and direct escape; forced first interaction when continuation is needed}\\
Interaction & Dust scattering versus absorption drawn probabilistically &
Implicit absorption; dust-albedo continuation & Weighted dust-scattering or absorber-yield continuation\\
Gas absorption & Immediate local combined-emissivity re-emission &
No immediate gas emission & Immediate emission for H$^0$, He$^0$, He$^+$; explicit remainder also emitted\\
Rates & Frequency-indexed path lengths, then frequency sums &
\both{Cross-section path integrals at each packet's actual energy; spectral bins retained separately}\\
Diffuse budget & Histories recycle launch normalization; extra sources possible &
Whole state emission launched separately & State emission split between history continuation and separate launch\\
Position & Absorption point; additional diffuse launches use cell centers &
Uniform inside emitting leaf & Absorption point for re-emission; explicit remainder uniform inside leaf\\
Low energies & Ordinary lines terminate; low-energy continuum travels if dust is present &
\both{Chosen source band can include FUV; general optical/IR line outputs are local emissivities, not full line transfer}\\
Thermal dust & Current dust emission CDF in absorption histories; temperature updated between passes &
\both{Gas-band grain absorption/scattering plus a separate optional dust-IR iteration}\\
Outer iteration & Transport, cell update, weighted grid convergence fractions; optional packet growth &
\both{Transport, rates, relaxed gas update, opacity update; fixed prescribed source count; final transport and gas update}\\
Conservation & Unchanged history normalization; photon number changes with frequency &
Separate source/state emission and band accounting & Gross band closure plus re-emission/net absorption; signed deposition and roulette\\
Diffuse noise & Ordinary sample count follows absorption histories &
Full-state explicit sample population & Remainder samples plus correlated descendants; no general noise advantage\\
Grid & Nested Cartesian grids; cell boundaries derived from axis midpoints &
\both{Octree leaves or dedicated Cartesian DDA; optional solution-driven refinement}\\
MPI & Source packet division; estimator/state reductions &
\both{Static cyclic or dynamic packet scheduling; distributed gas solves; node-shared state arrays}\\
Random histories & Intrinsic streams reseeded from clock/rank state in driver calls &
\both{QMC fixes launches by global index; interactions remain pseudorandom and rank/schedule dependent}\\
Observations & Escaped energy by origin, frequency, angle; auxiliary cube outputs &
Direct and dust-scattering observer estimates & Same, plus isotropic re-emission peel before roulette\\
Restart & Physical grid/state warm start &
Mapped gas/rate state, vacuum by stored $n_{\rm H}$; optional stopping histories & Mapped state as in B, reset stopping histories, first pass entirely explicit\\
\end{longtable}\endgroup

\section{What the three methods share}
\label{sec:same}
\Source These are steady radiation-transfer and gas-balance iterations,
not physical time integration of individual photons and gas parcels.
A pass fixes cell properties for transfer, accumulates radiation
estimators, then updates gas. In A, opacity and emission precede packets
and \code{updateCell} follows (\mc{iteration_mod.f90}{120--218,\allowbreak 295,\allowbreak 458--551,\allowbreak 828--833}).
In B/C, the loop calls \code{ionizing\_field\_and\_rates}, then the
ionization/thermal update and opacity rebuild
(\ms{main.f90}{200--237,\allowbreak 263}).

\Source All three sample spectra by cumulative energy-emission
probabilities, use straight flights through piecewise cell coefficients,
and estimate radiation from traversed material rather than requiring
absorption in a cell. A samples frequency indices and lengths
(\mc{photon_mod.f90}{730--756,\allowbreak 823--824,\allowbreak 1563--1574}); B/C invert an
energy-luminosity CDF and score exact-energy path integrals
(\ms{energy_sampler_mod.f90}{44--81,\allowbreak 203--254};
\ms{ion_score_mod.f90}{413--500}). This common principle does not imply
identical frequency discretization or variance.

\Inferred None of these packets is one physical photon. Even A's
``energy packet'' changes its represented photon number after a frequency
change. In steady state its normalization is an energy rate: A explicitly
describes \code{deltaE} in units of $10^{36}$ erg~s$^{-1}$
(\mc{photon_mod.f90}{272--279}); B/C use $W=wL_{\rm packet}$ and photon
rate $W/E$ (\ms{ion_score_mod.f90}{567--588}). Merely contrasting energy
and luminosity misses the substantive difference: how normalization
changes during a history.

\Source Both families include dust extinction/scattering, local emission
spectra depending on the current state, and MPI field reductions. A uses
scattering opacity divided by extinction and an HG direction when selected
(\mc{photon_mod.f90}{1672--1697,\allowbreak 1765--1797}); B/C use exact-energy dust
coefficients and HG scattering with weighted continuation
(\ms{ion_packet_mod.f90}{22--90}; \ms{raytrace_amr.f90}{402--460}).
These are common operations implemented with different Monte Carlo decisions.

\section{MOCASSIN versus the v1.00 leaf-emissivity algorithm}
\label{sec:ab}
\subsection{Packet normalization, spectra, and source sampling}
\Source A's packet stores frequency/index, position, direction, grid/cell
indices, and stellar/line flags, with no evolving weight
(\mc{common_mod.f90}{305--321}). Each star has $\delta E_j=L_j/N_j$;
an extended diffuse source can have a cell normalization
(\mc{set_input_mod.f90}{900--905}; \mc{photon_mod.f90}{65}).
Re-emission and escape retain that history's \code{deltaE(iStar)}.
``Equal energy'' is therefore a within-history statement here, not a
guarantee of identical normalization across every source/cell.

\Source B separates \code{energy\_eV}, bin label, \code{Lpacket}, and
\code{wgt}. Source launches set $w=1$ and $L_{\rm packet}=L_{\rm band}/N$
or the corresponding stratum value (\ms{define.f90}{131--162};
\ms{ion_band_mod.f90}{531--566,\allowbreak 2459--2466}). Components are drawn in
proportion to integrated band luminosity
(\ms{ion_band_mod.f90}{341--424,\allowbreak 595--623}); A loops over sources and
partitions prescribed counts among ranks (\mc{iteration_mod.f90}{474--496}).

\Source A selects an index with \code{getNu2}, then sets
$\nu=\code{nuArray(nuP)}$ (\mc{photon_mod.f90}{538--552,\allowbreak 730--756,\allowbreak 823--824}).
B integrates a nonnegative piecewise linear spectral luminosity density,
searches its independent CDF, and inverts within the sampler interval.
It then labels the energy with an output bin; it does not select that
output bin first and draw uniformly inside it
(\ms{energy_sampler_mod.f90}{44--81,\allowbreak 203--254};
\ms{ion_band_mod.f90}{2457--2466}). Continuous sampling refers to the
adopted representation, not an unapproximated Planck function. The CDF
is proportional to energy luminosity, not luminosity divided by photon energy.

A tabulated input spectrum reaches the two codes differently, and the same
file then gives different source spectra. For a spectrum file (a file name
as \code{contShape}, branch \code{'rauch'}), \Source A converts each input
point to $F_\nu$, assigns it to the frequency bin whose midpoint interval
contains it, and averages the points of a bin; a bin that receives no
input point keeps zero flux, and the bin probabilities are
$F_\nu\,\Delta\nu$ over their sum
(\mc{set_input_mod.f90}{351--357}; \mc{continuum_mod.f90}{216--283,\allowbreak 442--458}).
No interpolation fills the empty bins in this \code{rauch} branch,
which has the same binning expressions in the user's RT tree. The
\code{SB99} branch instead interpolates internal empty energy-flux bins;
in A it constructs photon flux before filling them, without recomputing
that flux afterward (\mc{continuum_mod.f90}{192--214}). A also initializes
\code{lamCount} outside the source loop, so multiple tabulated sources
reuse earlier counts; the RT tree moves that reset inside the loop and
recomputes SB99 photon flux after filling gaps
(\mc{continuum_mod.f90}{59,\allowbreak 92--99}).
\Source B integrates the piecewise linear spectrum through the file points. Where the frequency grid of A is
finer than the file, part of the band is therefore removed from A's
source, and the removed share varies with energy because A refines its grid
near the ionization edges.

\Measured Previously evaluated for one input (a 1220-point
stellar population spectrum, \code{nbins = 600}, $0.092$--$17$~Ryd): the
file has 42, 36 and 43 points in 13.6--24.6~eV, 24.6--54.4~eV and above,
against 51, 70 and 89 bins of A, of which 15, 34 and 46 are empty. At equal
hydrogen-ionizing photon rate, $Q({\rm He^0})/Q({\rm H^0})$ is 0.1297 for
the integrated file and 0.1054 for A's binned spectrum, a factor 1.230;
between 27 and 40~eV the integrated file carries 1.35 to 1.55 times the
photons of the binned one. Matching $Q({\rm H^0})$ between the codes thus
leaves the helium- and metal-ionizing parts of the source unmatched; an
identical source needs A's binned spectrum, zeros included, as the input of
B, or an input file resampled more finely than A's grid. These are quoted
source-spectrum reconstruction values, not a new MOCASSIN execution;
the output provenance is qualified in item 8 of Section~\ref{sec:unknown}.

\Source A uses intrinsic \code{random\_number}; isotropic directions
sample uniform cosine/azimuth, with separate plane/symmetry rules
(\mc{vector_mod.f90}{303--314}; \mc{photon_mod.f90}{560--675}). B/C offer
pseudorandom or scrambled Sobol launches. Spectrum, angles, and additional
position coordinates are indexed by global packet number on separate
stellar, diffuse, and FUV streams (\ms{qmc_mod.f90}{54--56,\allowbreak 78--177};
\ms{ion_band_mod.f90}{2427--2466,\allowbreak 2762--2777};
\ms{ionizing_field_mod.f90}{175--186,\allowbreak 429--450}). Arbitrary Sobol prefixes
are callable, but digital-net balance cannot be promised for every
non-power-of-two diffuse count.

\subsection{One flight: sampled absorption versus analytic attenuation}
\label{sec:flight}
\Source A draws $\tau_*=-\ln(1-u)$, accumulates cell optical depth,
and places the event at the corresponding distance inside the cell
(\mc{photon_mod.f90}{1177--1184,\allowbreak 1435,\allowbreak 1516--1527}). With dust it selects
scattering versus absorption using $\kappa_{\rm sca}/\kappa_{\rm ext}$.
Gas-only absorption selects diffuse emission; dust-only selects dust
emission; mixed cells select combined diffuse emission
(\mc{photon_mod.f90}{1672--1714,\allowbreak 1801--1812}). This connected decision
contains no separate H/He/metal-absorber selection.

\Source B first integrates attenuation along the complete chosen launch
ray. With physical segment length $\ell$, incident luminosity weight $W$,
and preceding optical depth $\tau_0$, its luminosity-length contribution is
\begin{equation}
 {\cal P}_{\rm direct}=W e^{-\tau_0}\frac{1-e^{-\kappa\ell}}{\kappa},
 \qquad {\cal P}_{\rm direct}\to W e^{-\tau_0}\ell\quad(\kappa\to0).
 \label{eq:direct}
\end{equation}
The code uses code lengths with unit conversion and stable
\code{expm1}/\code{log1p} expressions
(\ms{raytrace_amr.f90}{680--755,\allowbreak 803--869};
\ms{define.f90}{1556--1565}). Direct escape is $We^{-\tau_{\rm edge}}$.
Without dust scattering, B finishes this packet; gas absorption has
already contributed through the attenuated path.

\Source If scattering continues, B forces a first interaction:
\begin{equation}
 W_1=W(1-e^{-\tau_{\rm edge}}),\qquad
 \tau_1=-\ln[1-u(1-e^{-\tau_{\rm edge}})].
\end{equation}
That sampled incoming flight is excluded from duplicate path scoring.
B then keeps the scattering share by multiplying the weight by
$\kappa_{\rm sca}/\kappa_{\rm ext}$, chooses an HG direction, and follows
ordinary exponential flights with sampled lengths
(\ms{raytrace_amr.f90}{309--329,\allowbreak 402--465,\allowbreak 577--633}). Below $w=10^{-4}$,
roulette retains a packet with probability 0.1 and multiplies a survivor's
weight by ten. Absorbed gas power is not immediately converted into a
recombination descendant; it contributes to the gas solution supplying
separate diffuse launches.

\Inferred Analytic attenuation reduces randomness along a chosen ray;
it does not remove randomness in its direction, energy, or diffuse
position. A's probabilistic decisions and B's weighted survival are
different estimators. Neither alone establishes a numerical speed or
noise advantage.

\subsection{Radiation estimators and their actual consumers}
\label{sec:estimator}
\Source A accumulates both crossed and interaction-segment lengths,
schematically $\sum\delta E\ell/V$, in frequency-indexed arrays
(\mc{photon_mod.f90}{1563--1574,\allowbreak 1822--1833}). Stellar segments enter
\code{Jste}; nonstellar segments enter \code{Jdif} only with
\code{lgDebug}, otherwise also \code{Jste}. Thus ordinary \code{Jste}
is not purely stellar. Rate/heating consumers combine the separated
fields when debug separation is enabled
(\mc{update_mod.f90}{253--271,\allowbreak 1205--1248,\allowbreak 1492--1506}). The arrays have
bin-integrated normalization; conventional $J_\nu$ factors cannot be
assigned from their names alone (\mc{photon_mod.f90}{272--279};
\mc{iteration_mod.f90}{705--722}).

\Source B/C directly accumulate primary gas rates at the actual packet
energy. For H/He, in physical units the corresponding expressions are
\begin{equation}
 \Gamma_s=\frac1V\sum {\cal P}\frac{\sigma_s(E)}{E},\qquad
 H_s=\frac1V\sum {\cal P}\sigma_s(E)
             \max\!\left(1-\frac{E_{{\rm th},s}}E,0\right),
 \label{eq:rates}
\end{equation}
where $H_s$ heats one absorber, $E$ has units consistent with the
luminosity, and ${\cal P}$ is an attenuated direct or sampled
luminosity-length segment. The code scores these coefficients before
reduction and volume/unit normalization; gas-rate calculation calls the
reduced integrals rather than re-integrating output bins at representative
energies (\ms{ion_score_mod.f90}{413--500,\allowbreak 1039--1060};
\ms{gas_rates_mod.f90}{110--116}). Optional secondary-electron and
metastable-state integrals are additional terms. Metals likewise score
cached exact-energy rate/heating coefficients, rather than assuming
one threshold for every absorption contribution
(\ms{species_mod.f90}{525--538,\allowbreak 599--609}).

\Source Binned \code{jt\_ion} still receives path luminosity and has
physical consumers in dust radiation input and grain charge/heating,
as well as spectral output (\ms{ion_score_mod.f90}{435--441};
\ms{jtally_mod.f90}{198--226}; \ms{sedust_mod.f90}{259--274};
\ms{grain_photoelectric_mod.f90}{200--215}). Thus gas primary rates avoid
representative-bin energies, but the bins are not exclusively diagnostic.

\subsection{Diffuse emission, location, and nonionizing photons}
\label{sec:diffuseb}
\Source A forms continuum, line, and dust contributions from current
cell state. \code{recPDF} is a normalized gas/dust continuum CDF with selected
ionizing helium lines inserted into its frequency bins.
\code{totalLines} is the ordinary-line energy fraction, and
\code{linePDF} is a separate cumulative line distribution for debug
recording (\mc{emission_mod.f90}{113--159,\allowbreak 990--1129,\allowbreak 1179--1312}). A
replacement draws those probabilities while retaining history
normalization. In gas--dust cells its emission spectrum is combined,
not conditioned on an atomic absorber (\mc{photon_mod.f90}{907--990}).

\Source Ordinary absorption returns the actual location to the next
initialization (\mc{photon_mod.f90}{1524--1527,\allowbreak 2861--2869}). Additional
extended diffuse and resonance-related populations can start at cell
centers (\mc{photon_mod.f90}{208--213,\allowbreak 241--249,\allowbreak 882--903}).
Ordinary line packets are marked \code{lgLine}, optionally counted,
and terminate without normal trajectory transfer, even with dust.
A continuum below the hydrogen edge immediately escapes only without
dust; with dust it travels (\mc{photon_mod.f90}{370--482,\allowbreak 923--960}).
Resonance escape probabilities and associated dust-energy packets are an
additional approximation, not general tracking of every optical/IR line
(\mc{emission_mod.f90}{161--246,\allowbreak 3059--3099};
\mc{photon_mod.f90}{180--264}).

\Source B computes state luminosities $L_{c,i}$ for diffuse channels in
leaf $i$. It chooses the leaf by $\sum_cL_{c,i}$, position uniformly
inside it, isotropic direction, channel by energy luminosity, and energy
from that channel. For positive total luminosity,
\begin{equation}
 N_{\rm dif}=\min\left\{N,\,\max\left[1,\operatorname{nint}
 \left(\frac{L_{\rm dif}}{L_{{\rm packet,ion}}}\right)\right]\right\},
 \qquad L_{{\rm packet,dif}}=L_{\rm dif}/N_{\rm dif},
 \label{eq:ndif}
\end{equation}
with $N$ the source packet count (\code{no\_photons}). The limit acts only
where $L_{\rm dif}$ exceeds the ionizing source luminosity, as in a state
far more ionized than the source can keep (a fully ionized initial state);
each diffuse packet then carries more than a source packet and the
transported total stays $L_{\rm dif}$.
Zero luminosity launches no packets
(\ms{diffuse_mod.f90}{274--332,\allowbreak 399--424,\allowbreak 482--537,\allowbreak 677--763}). B tabulates
this after source transport, without a gas solve between populations;
both still see the frozen pass state
(\ms{ionizing_field_mod.f90}{155--186,\allowbreak 217--220}).

\Source B/C can extend the source band into FUV below the hydrogen edge,
with optional independent stratum normalization and count. Such photons
can affect dust, metals, and enabled metastable photoionization
(\ms{ion_band_mod.f90}{196--218,\allowbreak 531--566};
\ms{ionizing_field_mod.f90}{155--163};
\ms{ion_score_mod.f90}{437--488}). General line lists/leaf emissivities
are evaluated after the gas solution and summed locally, rather than
propagated as general line transfer (\ms{main.f90}{447--450};
\ms{lines_mod.f90}{45--80,\allowbreak 319--343}). Ionizing cascade lines carried in
diffuse channels are the explicit exception.

\subsection{Dust emission is organized differently}
\label{sec:dust}
\Source A includes dust continuum in the current combined emission CDF,
or a dust-only CDF without gas. Temperatures are updated after transfer
from grain/radiation energy balance
(\mc{emission_mod.f90}{1053--1129,\allowbreak 1342--1396};
\mc{photon_mod.f90}{994--1047}; \mc{update_mod.f90}{296--303,\allowbreak 1853--1925}).
Immediate re-emission uses that computed distribution; it is not an
identified event-by-event temperature correction from a difference of
Planck functions.

\Source B/C include grain absorption/scattering and exact-energy dust
heating in gas-band transfer. After the final gas pass they can compute
dust emission and invoke a separate dust-IR iteration. It samples leaf
spectra, uses luminosity-based systematic allocation to leaves, updates
absorbed power, and stops on its own relative-change tolerance or cap
(\ms{main.f90}{380--404};
\ms{dust_emis_transport_mod.f90}{244--299,\allowbreak 350--381,\allowbreak 533--589}). The shared
transport routine runs under a dust-IR flag; scoring sends absorption to
the dust estimator and returns before gas-rate scores
(\ms{ion_score_mod.f90}{425--433}). C does not immediately turn each
gas-band grain absorption into a dust-IR descendant in its r7 gas rule.

\subsection{Outer transport and gas updates}
\label{sec:iteration}
\Source A freezes opacity/emission, transports packets, reduces the
field, and updates local gas/dust states. B/C likewise update gas after
transport and rebuild opacity for the next pass
(\mc{iteration_mod.f90}{120--218,\allowbreak 295,\allowbreak 458--722,\allowbreak 828--889};
\ms{main.f90}{200--263}). B/C relax H/He fractions toward the local
solution through \code{ion\_relax}; their selected temperature is taken
directly (\ms{ion_balance_mod.f90}{1199--1222};
\ms{thermal_mod.f90}{338--358}). The separate outer and inner stopping
criteria, assigned defaults, packet-growth rules, and final-output
qualifications are compared in Section~\ref{sec:stopping}.

\subsection{Grids, MPI, observations, and restart}
\label{sec:implementation}
\Source A's Cartesian parent/child grids mark replaced cells with
negative child-grid indices. Transfer changes grids and resynchronizes
parent indices. Its connected 3D path computes distances to cell walls
from adjacent axis midpoints, selects the nearest wall, and updates
position and indices inline (\mc{grid_mod.f90}{835--875};
\mc{photon_mod.f90}{1206--1242,\allowbreak 1263--1376,\allowbreak 1394--1435,\allowbreak 1838--1878}).
The user's RT tree replaces this code with explicit-face DDA and separate
subgrid/multigrid routines; those added modules are not part of A.
A rejects \code{oneD}, despite retained spherical expressions
(\mc{set_input_mod.f90}{291--297}). B/C use octree leaves or dedicated
Cartesian DDA and can refine from the gas solution
(\ms{raytrace_amr.f90}{692--742,\allowbreak 764--956}; \ms{main.f90}{264--275}).
DDA is a traversal method, not a mesh hierarchy; the distinction here is
the actual boundary representation and connected calls.

\Source A divides packets by quotient/remainder and reduces fields
(\mc{iteration_mod.f90}{474--496,\allowbreak 616--722,\allowbreak 848--889}).
A reduces private arrays through temporary buffers; the user's RT tree
adds shared windows and revised reductions, without those additions
being part of this description of A. B/C assign global
indices cyclically or in dynamic batches; a dynamic dispatcher rank does
not transport. Gas solves are distributed and published state/rate
arrays node-shared, but not every tally is shared
(\ms{photon_schedule_mod.f90}{82--165,\allowbreak 189--247};
\ms{setup.f90}{75--81}; \ms{gas_state_mod.f90}{167--204};
\ms{gas_rates_mod.f90}{94--104}; \ms{jtally_mod.f90}{198--226}).

\Source A uses intrinsic pseudorandom streams. Each
\code{energyPacketDriver} call obtains the current seed, adds clock
milliseconds and \code{taskid}, and resets it
(\mc{photon_mod.f90}{68--87}). The fixed-seed input and one-time seeding
option belong to the user's RT tree, not A.
B/C Sobol launches can preserve their
global point set, but interactions/roulette/re-emission use stateful,
rank-associated pseudorandom draws (\ms{random_mt.f90}{947--957};
\ms{raytrace_amr.f90}{328,\allowbreak 385,\allowbreak 414--427,\allowbreak 464}). Fixed launches or seeds do
not guarantee bitwise outputs after rank/schedule/reduction changes.

\Source A bins escape by its current origin cell, frequency, and angle;
initialization resets that origin. Its driver writes escaped SED/cube
products (\mc{photon_mod.f90}{389--460,\allowbreak 679--681};
\mc{iteration_mod.f90}{738--741};
\mc{output_mod.f90}{2574--2581,\allowbreak 2647--2672}). A's nonsymmetric
azimuth bins use \code{atan(y/x)}, which loses the opposite-quadrant
distinction; the user's RT tree uses an \code{atan2}-based routine
(\mc{photon_mod.f90}{389--410}). B/C use observer-directed
peel at launches and scattering: they estimate contributions toward each
camera, trace attenuation, and use isotropic, inward-surface, or HG
weights (\ms{ionizing_field_mod.f90}{121--126,\allowbreak 185,\allowbreak 449};
\ms{ion_peel_mod.f90}{254--272,\allowbreak 292--356}). The identified A route does not
establish a matching MoCHII-style peel calculation.

\Source A's warm driver reloads state and reconstructs radiation setup;
\code{writeGrid} writes binary state
(\mc{mocassinWarm.f90}{39--59};
\mc{grid_mod.f90}{2646--2721,\allowbreak 2774--2816,\allowbreak 2843--2923,\allowbreak 3476--3488,\allowbreak 3546--3580}). B/C map
new leaf centers into old leaves and restore gas/available rate state,
with optional stopping histories in B
(\ms{previous_state_mod.f90}{204--242,\allowbreak 296--377};
\ms{main.f90}{103--119}). A stored leaf is vacuum for new gas where the
file's \code{n\_H} is zero; a file without \code{n\_H} (written before
2026-11-13) falls back to stored $n_e\le0$ with ionized fractions, and
stored neutral gas with $n_e\le0$ restarts at the $10^{-30}n_{\rm H}$
floor (\ms{previous_state_mod.f90}{213--235,\allowbreak 244--255}). These are physical warm starts, not established
in-flight packet/random-stream continuation.

\section{Convergence criteria and what their stopping flags mean}
\label{sec:stopping}

\Source The two families test different quantities and norms. Their
numerical tolerance values therefore cannot be ranked as if they were
tolerances on the same error. A local ionization/temperature solver, an
outer radiation/state iteration, and a final output update are also
different operations. The formulas below describe the inspected local
MOCASSIN tree and the shared MoCHII stopping code, not a new convergence
experiment. The v1.00 and v1.10 copies of
\code{fixed\_point\_distance\_mod}, \code{thermal\_mod}, and
\code{ion\_balance\_mod} are currently identical; C changes radiation
sampling and restart policy, not these mathematical stopping tests.

\subsection{MOCASSIN: a count of cells passing a relative-change test}
\label{sec:mocstop}

\Source For a gas-bearing cell the outer flag is
\begin{equation}
 q_i=\mathbf{1}\!\left[
 \left|\frac{x_{{\rm HI},i}^{\rm new}-x_{{\rm HI},i}^{\rm old}}
 {x_{{\rm HI},i}^{\rm old}}\right|\leq\epsilon_{\rm M}\right],
 \qquad \epsilon_{\rm M}=\code{XHILimit}.
 \label{eq:moccell}
\end{equation}
This uses the old neutral fraction as denominator, not the absolute
change in the ionized fraction. Temperature, electron density, and
other ion fractions are not additional outer gas-cell requirements.
In gas with dust, the same H I flag is used; only in dust without gas
is the tested relative change instead that of the size/species-weighted
mean \code{Tdust(0,0,cellP)}
(\mc{update_mod.f90}{159--169,\allowbreak 319--330,\allowbreak 513--521,\allowbreak 1939--1946}). The literal
relative expressions do not protect a zero old denominator at their
point of use.

\Source For grid $g$, let $\mathcal A_g$ be its counted active
cells (restricted by light-echo settings), $m_g=\code{grid(g)\%nCells}$,
and $c_g=\sum_{i\in\mathcal A_g}q_i/|\mathcal A_g|$, with $c_g=0$ for an
empty counted set. The original driver computes
\begin{equation}
 C_n=100\frac{\sum_g m_gc_g}{\sum_g m_g},\qquad
 \hbox{outer success if }C_n\geq\code{minConvergence}.
 \label{eq:mocpercent}
\end{equation}
This is a mean of grid passing fractions weighted by \code{nCells},
not generally a direct sum over the counted cells, and uses no physical
volume, mass, or luminosity weights. For one ordinary grid it reduces
to the passing-cell percentage. The user's RT tree instead sums the
converged and counted-active cell totals across grids. A cell receiving no sampled radiation
keeps its old state but its convergence flag remains zero; it still
belongs to the ordinary active-cell denominator. Thus unchanged because
of no hits is not accepted as convergence
(\mc{iteration_mod.f90}{807--816,\allowbreak 955--998,\allowbreak 1056--1061};
\mc{update_mod.f90}{92--156}). The code can declare the percentage test
satisfied even when some local thermal/ionization solves set failure
counters: those counters and \code{lgBlack} do not veto the H I flag
(\mc{update_mod.f90}{493--520}).

\Source The actual input records and fresh-input defaults are:
\begin{center}\small
\setlength{\tabcolsep}{4pt}
\begin{tabular}{@{}P{.32\linewidth}P{.33\linewidth}P{.29\linewidth}@{}}
\toprule Input record & Internal controls & Assigned default\\\midrule
\code{convLimit value} & \code{XHILimit} & 0.05\\
\code{maxIterateMC count percent} & \code{maxIterateMC}, \code{minConvergence} & 30, 95\\
\code{autoPackets increment factor maximum} & \code{convIncPercent}, \code{nPhotIncrease}, \code{maxPhotons} & Disabled; numeric fields 0, 0, 0\\
\code{nPhotons count} & Initial stellar packet budget & No assigned default established in the input block\\
\code{quantumHeatGrain radius percent} & Enables qHeat; \code{minaQHeat}, \code{minConvQHeat} & Disabled; $10^{-3}$ $\mu$m, 99\\
\code{quantumHeatGrainParameters temperature bins writeFlag} & \code{Tmax}, \code{nTbins}, \code{lgWritePss} & 700 K, 300, false\\\bottomrule
\end{tabular}
\end{center}
These are input keywords, not interchangeable aliases for internal
variable names. In particular, \code{XHILimit} and \code{minConvergence}
are not standalone input keys. Parser assignments and reads are at
\mc{set_input_mod.f90}{34,\allowbreak 67--89,\allowbreak 105--117,\allowbreak 191--199,\allowbreak 281--291,\allowbreak 377--385,\allowbreak 413--416};
state declarations are at \mc{common_mod.f90}{33,\allowbreak 76--93,\allowbreak 402--443,\allowbreak 460}.
Warm-state loading can restore controls from its records, so parser
defaults are not assertions about every warm run
(\mc{grid_mod.f90}{3097,\allowbreak 3107}).

\Source After a completed pass, automatic packet growth requires all of
\begin{equation}
 n>1,\quad C_n<95,\quad C_{n-1}>0,\quad
 \code{lgAutoPackets},\quad \widetilde N_n<\code{maxPhotons},\quad
 \frac{C_n-C_{n-1}}{C_{n-1}}\leq\code{convIncPercent}.
 \label{eq:mocgrowth}
\end{equation}
The implemented maximum comparison uses
$\widetilde N_n=N_{1,n}+\sum_{j=1}^{n_\star}N_{j,n}$, because
\code{nPhotonsTot} starts at \code{nPhotons(1)} before adding all stars.
The first star is counted twice; extended-diffuse packets are excluded.
The user's RT tree starts the sum at zero and removes this double count.
The value 95 here is literal, independent of a changed
\code{minConvergence}. The increment is a dimensionless relative change,
not percentage points. When triggered, each stellar count is rounded
from $N_j\code{nPhotIncrease}$ and its normalization divided by that
factor; an enabled extended-diffuse count is also increased. The maximum
is checked \emph{before} multiplication, without a clamp afterward, so
it is not a strict upper bound on the next count. A previous zero
percentage prevents growth. This logic runs before testing either the
convergence percentage or \code{maxIterateMC}; the last pass can therefore
change counts without executing another, enlarged population
(\mc{iteration_mod.f90}{1088--1112,\allowbreak 1129--1160}).

\Source There is no eight-sample minimum. On either percentage success
or exhaustion of the outer cap, A writes grid/gas results and returns
without another recursive pass; cap exhaustion has a separate message but does not
provide a strict abort analogous to MoCHII's
\code{require\_convergence}. No additional transport of the final updated
gas is invoked on these exit branches
(\mc{iteration_mod.f90}{828--889,\allowbreak 1135--1160}). Exiting normally is
therefore not itself evidence that the configured percentage was reached.

\subsection{MoCHII: maximum steps, normalized integrated histories, or MOCASSIN's count}
\label{sec:mochstop}

\Source Write $\Delta x=x^{\rm new}-x^{\rm old}$, including the
configured outer H/He relaxation. For positive gas volume, the measures
use $D(Q)=\max(\sum_i V_i Q_i^{\rm new},\sum_i V_i Q_i^{\rm old})$
and are
\begin{align}
 d_{\rm H}&=\frac{\sum_i V_i|\Delta x_{{\rm HII},i}|}
                   {D(x_{\rm HII})},\\
 d_{{\rm He},1}&=\frac{\sum_i V_i|\Delta x_{{\rm HeI},i}|}
 {D(x_{\rm HeII}+x_{\rm HeIII})},\\
 d_{{\rm He},3}&=\frac{\sum_i V_i|\Delta x_{{\rm HeIII},i}|}
 {D(x_{\rm HeIII})},\qquad
 d_T=\frac{\sum_i V_i n_{e,i}^{\rm new}|\Delta T_i|}
 {D(n_e T)}.
 \label{eq:integratedsteps}
\end{align}
Tiny denominator floors are applied where the code requires them. The
He III measure is active only when its integrated abundance is strictly
greater than 1\% of integrated ionized helium, rather than imposing that
cut independently in each leaf
(\ms{ion_balance_mod.f90}{1225--1254,\allowbreak 1340--1349};
\ms{thermal_mod.f90}{368--390,\allowbreak 436--447}). These are normalized integrated
absolute steps: opposite signs do not cancel in their numerators.

\Source With \code{conv\_crit='cell'}, the outer test is
\begin{equation}
 \max_i|\Delta x_{{\rm HII},i}|<\code{gas\_tol},\qquad
 \max_i\frac{|\Delta T_i|}{\max(T_i^{\rm old},T_i^{\rm new})}<\code{gas\_tol\_te}
 \quad\hbox{when temperature is solved}.
 \label{eq:cellstop}
\end{equation}
It does not impose a helium test or an eight-iteration window.
With \code{conv\_crit='vol'} and optional floor stopping disabled
(the default), each active composition measure instead
passes the history test below, and temperature still requires the
latest $d_T<\code{gas\_tol\_te}$
(\ms{main.f90}{267,\allowbreak 566--591,\allowbreak 621--635}). The hydrogen balance residual and
changes in every metal stage are not extra conditions in these outer
logical expressions (\ms{main.f90}{566--591}).

\Source The optional \code{conv\_stop\_on\_floor} set to true applies to
\code{vol} in both MoCHII versions.  Each active composition measure
and solved temperature must pass either eight-update extrapolation
($\rho<0.95$, fitted remaining change below tolerance), or a floor test
($\rho\ge0.95$, latest step below ten times tolerance, and $t<2.365$
in each of three consecutive nonoverlapping eight-update signed windows).
The temperature numerator is $\sum Vn_e^{\rm new}\Delta T_e$, divided
by the same old/new maximum as $d_T$.  This is a statistical stop under
a stationary-noise assumption, not proof of equilibrium or exclusion
of drift below the finite-window resolution.  Any floor acceptance
sets \code{CONVERGD=false}, while \code{require\_convergence} accepts
the selected stop.  Temperature history is not saved: even a restart
restoring composition requires eight fresh temperature updates for
extrapolation and 24 for its floor test.  Re-refinement clears all
histories.  The option defaults to false; the quoted comparisons retain
their stated stopping settings.  See the convergence section of
\path{docs/MoCHII_UserGuide.pdf} and \path{src/fixed_point_distance_mod.f90}.

\Source With \code{conv\_crit='cell\_fraction'}, MoCHII applies
Eq.~\eqref{eq:moccell} leaf by leaf with
$\epsilon_{\rm M}=\code{conv\_xHI\_tol}$ (default 0.05) and stops when the
number of passing leaves divided by the number of leaves with gas is at
least \code{conv\_cell\_fraction} (default 0.95), the one-grid form of
Eq.~\eqref{eq:mocpercent}: a count with no volume weight, also on an AMR
grid. Temperature, \code{gas\_tol}, \code{gas\_tol\_te}, and the drift test
do not enter it; \code{gas\_niter} remains the cap. A leaf that no packet
crossed in the iteration fails and stays in the denominator, as a MOCASSIN
cell without a hit. A crossing is a positive path tally $\sum\ell L$ over
every segment scored in the leaf, of every band the gas iteration
transports: the ionizing band and, with \code{add\_fuv}, the FUV band
from \code{efuv\_min} to the Lyman limit
(\ms{main.f90}{261--267,\allowbreak 566--591};
\ms{ion_balance_mod.f90}{1266--1323};
\ms{ion_score_mod.f90}{438,\allowbreak 1017--1020}). The share and the share
of leaves without a crossing are printed every iteration whatever
\code{conv\_crit} is. This reproduces A's leaf test and count; in this
mode the start state also follows A's (Section~\ref{sec:cfdiff}, item 9).
Section~\ref{sec:cfdiff} lists what is identical and everything that
still differs.

\Source The declaration default for \code{conv\_crit} is \code{cell}.
The composition and thermal tolerances are each 0.001; outer
\code{ion\_relax} defaults to 1. Drift testing, strict convergence,
floor stopping, and restart-history restoration default to false
(\code{conv\_drift\_test}, \code{require\_convergence},
\code{conv\_stop\_on\_floor}, and
\code{restart\_iteration\_history}, respectively;
\ms{define.f90}{937--966,\allowbreak 1038,\allowbreak 1118--1122,\allowbreak 1281--1283}). The default
\code{gas\_niter=0} means one fixed-state transfer, not an attempted
iterative gas solve. A positive value sets the outer cap; there is no
separate default 30-pass solve. Setup restricts the criterion name and
relaxation range, and notes that a drift flag outside \code{vol} has no
effect; enabling floor stopping outside \code{vol} is refused
(\ms{setup.f90}{154--176,\allowbreak 694--697}).

\Source With optional floor stopping disabled, \code{vol} requires eight samples in every active composition
history. For positive $d_k$ in that window, fit
\begin{equation}
 b=\frac{\sum(k-\bar k)(\ln d_k-\overline{\ln d})}
          {\sum(k-\bar k)^2},\qquad
 \rho=e^b,\qquad d_{\rm fit}=e^{\overline{\ln d}+b(m-\bar k)},
 \quad m=8.
\end{equation}
Thus $\rho$ is the exponential of the fitted slope, not the slope itself.
When $\rho<0.95$, the available distance estimate is
\begin{equation}
 e=\frac{d_{\rm fit}\rho}{1-\rho},\qquad
 \hbox{accept if }e<\code{gas\_tol}.
 \label{eq:distancecriterion}
\end{equation}
It is a geometric-tail estimate, not a measured norm of the exact
solution error. A nonpositive step, including a zero, or $\rho\geq0.95$
makes this estimate unavailable. After the independent eight-sample
requirement, the unavailable-estimate branch instead requires
\begin{equation}
 d_{
\rm latest}<\code{gas\_tol},\qquad
 t_s=\frac{|\bar s|}{\operatorname{sd}(s)/\sqrt8}<2.365
 \quad\hbox{only if \code{conv\_drift\_test} is true}.
 \label{eq:floorcriterion}
\end{equation}
Signed $s$ uses the same normalization as $d$ with $\Delta x_{\rm HII}$,
$-\Delta x_{\rm HeI}$, or $\Delta x_{\rm HeIII}$ in the numerator.
Zero signed mean gives $t_s=0$; nonzero mean with zero scatter gives an
effectively infinite statistic. The drift test is applied in this
fallback branch, not additionally to an available decaying-tail
estimate (\ms{fixed_point_distance_mod.f90}{80--89,\allowbreak 134--246};
\ms{ion_balance_mod.f90}{1363--1373}).

\Inferred Calling the second branch a Monte Carlo floor does not prove
that noise, rather than slow physical change, caused it. The drift
statistic can detect a coherent signed change that small steps alone
miss. Its nominal Student threshold does not establish a universal
confidence guarantee for serially correlated radiation/state iterates.
Nor does it allow arbitrarily large stationary noise: the positive
absolute step must still lie below \code{gas\_tol}. MoCHII does not
increase the prescribed source count automatically in this gas loop
(\ms{setup.f90}{70}; \ms{main.f90}{220--337}); varying diffuse counts
are a different mechanism, given by Equation~\eqref{eq:ndif}.

\subsection{Inner gas and dust solves are separate certificates}
\label{sec:innerstop}

\Source A's executable \code{maxIterateT=50} limits thermal evaluations;
a commented-out value 25 is not the active limit. The user's RT tree
uses 25.
The heated-gas test is $|H-C|<0.01H$, a heating/cooling residual rather
than a temperature step. A bracket is sought with factor-1.2 changes
and then residual interpolation. At the cap, a bound or bracket midpoint
is retained, failure diagnostics are set, and execution continues;
ionization is not reevaluated at that final replacement temperature
(\mc{update_mod.f90}{49,\allowbreak 81--89,\allowbreak 360--370,\allowbreak 396--520}). The ionization
recursion has \code{maxIterateX=25}; it tests relative H I, He I, and
He II changes against 0.01, allowing equality. Electron density is
recomputed from ionic charge, without an independent density residual
or density-change tolerance. Exceeding those inner tests at the cap
returns the last values with failure flags; the outer H I percentage
can nevertheless pass (\mc{update_mod.f90}{1342,\allowbreak 1568--1627}). Grain charge
has its own 100-call cap when that calculation is used
(\mc{update_mod.f90}{81--89,\allowbreak 590--599}). These are compiled limits, not
\code{maxIterateMC} input values.

\Source MoCHII's thermal solve evaluates heating minus cooling at
\code{te\_min}/\code{te\_max}, default 1000/50000 K. If no bracket lies
inside, it takes the corresponding bound. Otherwise up to 60 geometric
bisections stop at $T_{\rm hi}/T_{\rm lo}-1<10^{-5}$, followed by a
fresh ionization evaluation at the selected midpoint. There is no
separate normalized heat/cool residual acceptance test or post-cap
failure flag for this bisection
(\ms{define.f90}{1280--1283}; \ms{thermal_mod.f90}{299--324}). A small
outer temperature step at a clamped bound does not certify an
unconstrained thermal root.

\Source Its ion/electron-density loop allows 200 iterations, mixing the
new charge density halfway toward the old value with a $10^{-30}n_H$
floor. It requires $|\Delta n_e|\leq10^{-10}n_e$ and, for enabled coupled
processes, the corresponding small hydrogen/helium-fraction changes.
The optional subsequent metal charge-neutrality update also uses at
most 200 half-damped iterations with the density criterion
(\ms{ion_balance_mod.f90}{266,\allowbreak 416--452,\allowbreak 459--465,\allowbreak 971--989}). Cap exhaustion
returns the last iterate without a dedicated warning/abort or failure
flag. The separate hydrogen balance root has its own residual and
bracket limits (\ms{ion_balance_mod.f90}{838--902}). None is the same as
\code{gas\_tol}, and there is no MoCHII input named \code{maxinneriter}
in the inspected source.

\Source Dust convergence must also be identified by calculation. A's
dust-only outer flag tests the relative weighted mean dust temperature,
not every size temperature. Its optional qHeat CDF construction is
activated only for enabled grains with radius below \code{minaQHeat},
\code{convPercent>minConvQHeat}, and outer pass number greater than one.
The scalar \code{convPercent} is last assigned while visiting each grid
at the end of the preceding pass; it is the last visited grid's
percentage, not necessarily the global \code{totPercent}. qHeat itself
forms a finite temperature/enthalpy-grid recurrence; \code{nTbins=300}
and \code{Tmax=700} K are default discretization controls, not a residual
iteration cap/tolerance. There is no separate qHeat success flag or
required extra outer pass. With default outer target 95\% and qHeat
activation strictly above 99\%, enabling qHeat does not guarantee it
runs before ordinary stopping
(\mc{common_mod.f90}{33,\allowbreak 76--93};
\mc{set_input_mod.f90}{79--83,\allowbreak 191--199};
\mc{iteration_mod.f90}{295,\allowbreak 955--998,\allowbreak 1051--1061,\allowbreak 1135--1160};
\mc{emission_mod.f90}{1071--1078,\allowbreak 1350--1358,\allowbreak 1638--1669}).

\Source MoCHII's optional dust-IR loop is separate from its gas stop.
It tests integrated dust-emission luminosity change
$|L_{\rm new}-L_{\rm previous}|/\max(L_{\rm new},\mathrm{tiny})$
against \code{dust\_tol}, default 0.001, for at most
\code{max(dust\_niter,1)} passes, with declared \code{dust\_niter=2}.
A remaining change at the cap produces a warning and continues; this is
not a proof of stable temperature distributions in every leaf or of the
gas criterion. Algebraic direct-split energy closure is checked
separately from stochastic path-plus-escape closure
(\ms{define.f90}{528--538};
\ms{dust_emis_transport_mod.f90}{249,\allowbreak 311--337,\allowbreak 374--400}).

\subsection{Stops, final output, restart, and grid changes}
\label{sec:stopoutputs}

\Source With a positive gas iteration budget, MoCHII first records the
outer result in \code{CONVERGD} and its pass count. Exhausting the cap
warns; \code{require\_convergence=.true.} aborts before normal final
output through \code{abort\_run} (an \code{ERROR [rank 0]:} line, then
\code{MPI\_ABORT} with exit status~1), while false allows the result with
its failed flag (\ms{main.f90}{341--381}; \ms{utility.f90}{396--399}). Missing restart metal-rate state can force strict reconvergence, an
exception to the permissive fresh-run default
(\ms{previous_state_mod.f90}{353--360}). A permitted
gas-solving run then performs an additional final transfer, or imaging
transfer, and another gas update. That update is still subject to
\code{ion\_relax}; \code{FINALDX}/\code{FINALDTE} change, but the outer
\code{CONVERGD} value is not retested. \code{STATESOL} reports an update
in the written field, not an independently verified equilibrium residual
or transport/state fixed point
(\ms{main.f90}{396--420}; \ms{gas_rates_mod.f90}{243--248}). For
\code{gas\_niter<1}, no final gas solve is requested. This differs from
A's direct output/exit after its last ordinary cell update.

\Source Re-refinement forces MoCHII to continue and clears all
composition and temperature histories, even if an old-grid criterion had passed. It
also invalidates raw H/He field estimates and defers an iterate
checkpoint until transfer has taken place on the new grid
(\ms{main.f90}{287--304}; \ms{amr_refine_mod.f90}{223--228};
\ms{ion_score_mod.f90}{340--343}). Ordinary B restart can optionally
restore composition histories; by default it does not. Temperature
history always starts empty in both modes. C always clears the
composition histories too,
regardless of the requested restoration flag. Initial estimator setup
makes its first transfer entirely explicit; a later valid field permits
absorber re-emission. Accordingly a fresh \code{vol} history requires
eight new samples, whereas a B history restored from a checkpoint can
already contain earlier samples
(\ms{previous_state_mod.f90}{311--323,\allowbreak 364--392};
\ms{main.f90}{115--119,\allowbreak 216--223}; \ms{diffuse_mod.f90}{339--347}). This is a changed
warm-start sequence, not a changed numerical tolerance. C adds no new
outer acceptance test for every emission channel or the distribution
of large history weights.

\clearpage
\begingroup\small
\setlength{\tabcolsep}{4pt}\renewcommand{\arraystretch}{1.12}
\begin{longtable}{@{}P{.15\linewidth}P{.27\linewidth}P{.26\linewidth}P{.26\linewidth}@{}}
\caption{Stopping-rule comparison with optional floor stopping disabled. Defaults are assigned fresh-input
values, not universal recommended tolerances. Source evidence and
exceptions are given in Section~\ref{sec:stopping}.}\label{tab:stops}\\
\toprule Aspect & MOCASSIN (A) & v1.00 algorithm (B) & v1.10 re-emission (C)\\\midrule\endfirsthead
\toprule Aspect & MOCASSIN (A) & v1.00 algorithm (B) & v1.10 re-emission (C)\\\midrule\endhead
\bottomrule\endfoot
Gas quantity & Relative H I change, $\leq0.05$ by default &
\both{Default cell: maximum absolute H II step; optional vol: normalized integrated H/He histories; optional cell\_fraction: A's relative H I change}\\
Spatial weight & \code{nCells}-weighted mean of grid passing fractions; default target 95\% &
\both{Cell maximum, or volume-integrated composition; thermal vol step weights $Vn_e$; cell\_fraction: count of passing leaves over leaves with gas, default target 95\%}\\
Outer thermal test & No extra gas-temperature requirement; dust-only uses mean grain temperature &
\both{Maximum relative temperature step for cell; latest integrated $d_T$ for vol; none for cell\_fraction}\\
Tolerances & \code{convLimit} sets 0.05; second \code{maxIterateMC} value sets 95\% &
\both{\code{gas\_tol=0.001}, \code{gas\_tol\_te=0.001}; different norms from A; cell\_fraction: \code{conv\_xHI\_tol=0.05}, \code{conv\_cell\_fraction=0.95}}\\
Noise & Optional conditional packet growth; no signed-drift test &
\both{Fixed prescribed source budget. Vol floor still needs step below tolerance; optional signed-drift test defaults off}\\
Minimum & No enforced multi-pass history minimum &
Cell can pass first update; vol requires 8 samples, possibly restored & Same tests; restart discards old samples\\
Maximum & \code{maxIterateMC=30} &
\both{\code{gas\_niter}; default 0 is one fixed-state transfer; positive input caps the gas loop}\\
At cap & Writes results and exits with cap message &
\both{Warns and may output; \code{require\_convergence=true} aborts before normal output; default false}\\
Inner gas T & Heat/cool residual below 1\%; at most 50 evaluations &
\both{Bracket ratio below $10^{-5}$ or 60 bisections; bound clamping possible}\\
Inner ions/$n_e$ & H I/He I/He II relative changes at most 1\%; 25 evaluations; no independent $n_e$ test &
\both{200 iterations with $n_e$ relative step $10^{-10}$, plus enabled coupled-ion tests; optional second charge loop}\\
Inner failure & Diagnoses failed local solves without vetoing outer H I flag &
\both{These fixed-cap loops can return last iterate without an explicit failure flag; outer stop remains a state-change test}\\
Dust/qHeat & Mean dust-temperature change in dust-only mode; qHeat gated but has no separate stopping certificate &
\both{Optional separate dust-IR luminosity-change loop; default 2 passes and tolerance 0.001}\\
After success & Output; no added final MC transfer on the exit branch &
\both{Gas-solving run: final transfer/imaging and gas update; outer flag is not retested; inspect final-step diagnostics}\\
Grid/restart & Physical warm state can restore controls; no 8-step norm here &
Refinement clears histories; restart optionally restores them & Same refinement reset; restart always clears histories and begins with explicit emission\\
\end{longtable}\endgroup
\FloatBarrier

\subsection{\code{cell\_fraction} and A's test: what is identical and what still differs}
\label{sec:cfdiff}

\Source The list below was checked line by line against source A and the
current v1.10 source on November 15, 2026; the \textsf{S:} line numbers
here are those of that source, after the start-state rule of item 9 was
added. B and C share every \MoCHII\ line cited here, except where item 3
says otherwise. The compared inputs are the clumpy model of
\path{clumpy_test/}: A's \path{mocassin/input/input.in} (one grid,
$129^3$ cells) and the \MoCHII\ inputs \path{mochii_clmp1_cell_fraction/}
and \path{mochii_clmp2_cell_fraction/}.

\begingroup\small
\setlength{\tabcolsep}{4pt}\renewcommand{\arraystretch}{1.12}
\begin{longtable}{@{}P{.17\linewidth}P{.37\linewidth}P{.40\linewidth}@{}}
\caption{A's convergence test and \code{conv\_crit='cell\_fraction'}:
the identical parts.}\label{tab:cfsame}\\
\toprule Aspect & MOCASSIN (A) & \MoCHII\ \code{cell\_fraction}\\\midrule\endfirsthead
\toprule Aspect & MOCASSIN (A) & \MoCHII\ \code{cell\_fraction}\\\midrule\endhead
\bottomrule\endfoot
Leaf test & $|x^{\rm new}_{\rm HI}-x^{\rm old}_{\rm HI}|/x^{\rm old}_{\rm HI}\leq\code{XHILimit}$, set by \code{convLimit} (0.05) (\mc{update_mod.f90}{513--521}) & The same with \code{conv\_xHI\_tol} (0.05) (\ms{ion_balance_mod.f90}{1306--1323}); $x^{\rm new}$ is taken after \code{ion\_relax}, default 1 and not set in the clumpy inputs, so the same there; a zero old fraction is guarded, where A divides by zero\\
Count & Converged cells over active cells; a mean of grid shares weighted by \code{nCells} over several grids, a plain count on the one grid of the clumpy input (\mc{iteration_mod.f90}{955--998,\allowbreak 1056--1061}) & Converged leaves over leaves with gas, no volume weight, also on an AMR grid (\ms{main.f90}{526--535}); the several-grid weighting has no counterpart\\
Uncrossed cell & Not converged and kept in the denominator: \code{lgHit} needs a positive \code{Jste}, which also holds the diffuse packets; with \code{lgDebug} they go to \code{Jdif} instead, and a positive \code{Jdif} then counts as a hit as well (\mc{update_mod.f90}{106--125,\allowbreak 137--157}; \mc{photon_mod.f90}{1562--1574}); flags zeroed before every pass (\mc{iteration_mod.f90}{807--816}) & Not converged and kept in the denominator; a crossing is any scored path segment of any transported band (\ms{ion_balance_mod.f90}{1312--1316})\\
Temperature & Not tested & Not tested (\ms{main.f90}{518--519})\\
Stop & Share $\geq$ \code{minConvergence} (95) (\mc{iteration_mod.f90}{1135}) & Share $\geq$ \code{conv\_cell\_fraction} (0.95) (\ms{main.f90}{519})\\
\end{longtable}\endgroup

\begingroup\small
\setlength{\tabcolsep}{4pt}\renewcommand{\arraystretch}{1.12}
\begin{longtable}{@{}P{.03\linewidth}P{.29\linewidth}P{.33\linewidth}P{.25\linewidth}@{}}
\caption{A's convergence test and \code{conv\_crit='cell\_fraction'}:
what still differs, and whether its effect was measured.}\label{tab:cfdiff}\\
\toprule & MOCASSIN (A) & \MoCHII\ \code{cell\_fraction} & Effect\\\midrule\endfirsthead
\toprule & MOCASSIN (A) & \MoCHII\ \code{cell\_fraction} & Effect\\\midrule\endhead
\bottomrule\endfoot
1 & \code{autoPackets 0.20 2. 210000000} in the clumpy input multiplies \code{nPhotons} by 2 when the relative gain of the share is $\leq0.20$, the share is below 95, and the total is below the maximum (Eq.~\ref{eq:mocgrowth}; \mc{iteration_mod.f90}{1088--1112}). The first star is counted twice in \code{nPhotonsTot}, so A's run of $10^8$ packets kept $2\times10^8$ after one increase, where the RT tree went to $4\times10^8$ & Fixed \code{no\_photons} (and \code{no\_photons\_fuv}) in every iteration. The share can stall below the threshold on the Monte Carlo noise of $x_{\rm HI}$ & Stall measured: 80.42\% after 20 iterations at $10^6$ packets in a fixed-$T_e$ Str\"omgren sphere (\path{scratch_v1.10_claude/cell_fraction_check/v1.10/strom_cf}). The effect of a packet increase on \MoCHII\ was not measured\\
2 & Writes the state of the iteration that passed (\mc{iteration_mod.f90}{1135--1146}) & One more transport and one more solve after the loop (\ms{main.f90}{338--379}): the written state is one update past the tested one; \code{FINALDX}/\code{FINALDTE} record that step & Not measured\\
3 & Re-emits at the absorption point within the pass & B: leaf emission from the state of the previous solve; C: re-emission at the absorption point (Sections~\ref{sec:diffuseb} and \ref{sec:ac}) & Not measured under \code{cell\_fraction}\\
4 & Sampled absorption; $10^8$ packets in iterations 1--5 and $2\times10^8$ in 6--7 (\path{mocassin/output/summary.out}), $2\times10^8$ then $4\times10^8$ in the RT tree & Analytic attenuation along the path (Section~\ref{sec:flight}); $10^7$ packets in each iteration in clmp2, $10^8$ in clmp1, plus \code{no\_photons}/4 FUV packets & The noise of $x_{\rm HI}$ at each count was not measured in either code\\
5 & Ion balance stops at relative changes of H$^0$, He$^0$, He$^+$ $\leq0.01$ or at \code{maxIterateX}=25 (\mc{update_mod.f90}{83--84,\allowbreak 1342,\allowbreak 1602--1623}); thermal residual $|H-C|<0.01H$ or \code{maxIterateT}=50 (\mc{update_mod.f90}{49,\allowbreak 88--89,\allowbreak 404}). A cell at a cap is flagged (\code{lgBlack}, \mc{update_mod.f90}{495,\allowbreak 508--511}) and still enters the count & Section~\ref{sec:innerstop} & Not measured\\
6 & No temperature bounds; only $T\leq0$ is reset to 1~K (\mc{update_mod.f90}{460}) & \code{te\_min}, \code{te\_max} (\ms{define.f90}{1273--1274}; $10^3$ and $3\times10^4$~K in the clumpy inputs) & Not measured\\
7 & Packets from \code{nuMin} = 0.092~Ryd = 1.25~eV in the clumpy input; a packet stops where it is absorbed & No packet below \code{efuv\_min} = 4.78~eV (\code{eion\_min} without \code{add\_fuv}); the analytic edge walk scores leaves beyond the point a sampled packet would be absorbed & No-packet share 0.00\% in all 12 iterations of the two clumpy runs; the effect on the converged share was not measured\\
8 & \code{maxIterateMC 30 95.} (30 iterations) & \code{gas\_niter} = 40 in the clumpy inputs & None while a run stops before 30\\
9 & Start (\mc{grid_mod.f90}{1542,\allowbreak 1564--1601}): $X({\rm H^0})$ = \code{H0in} = $10^{-5}$, $X({\rm H^+}) = 1-X({\rm H^0})$; $X({\rm He^0}) = X({\rm H^0})$, $X({\rm He^+}) = 1-X({\rm He^0})$, $X({\rm He^{2+}}) = 0$; every metal: stage I $= X({\rm H^0})$, higher stages 0; $n_e = n_{\rm H}$; $T_e$ = \code{TeStart} (8000~K in the clumpy input) & Namelist start in this mode: $x_{\rm HI}$ = \code{xHI\_init} as the input gives it (not fixed at $10^{-5}$); helium as hydrogen, $x_{\rm HeI} = x_{\rm HI}$, $x_{\rm HeII} = 1-x_{\rm HI}$, $x_{\rm HeIII} = 0$, unless the input gives \code{xHeI\_init} or \code{xHeII\_init} (\ms{setup.f90}{1029--1032}; \ms{gas_state_mod.f90}{86--168}); until the first solve every metal at stage I $= x_{\rm HI}$, higher stages 0, not normalized, as in A (\ms{main.f90}{111--119}; \ms{species_mod.f90}{535--562}), applied once, so a re-refinement keeps the evolved cache, after which every solve forms the stages from the rates (\ms{species_mod.f90}{1902--1933}); $T_e$ = \code{te\_start} (the thermal solve; \code{te\_fixed} is the held $T_e$ of \code{solve\_te = F}). Still different: $n_e$ starts from the H and He fractions, $1.1\,n_{\rm H}$ at $y_{\rm He}=0.1$, where A takes $n_{\rm H}$ & Not yet measured: the clumpy runs quoted below started from $x_{\rm HI}=10^{-3}$, neutral helium and $T_e=10^4$~K\\
10 & A relative $x_{\rm HI}$ test can pass prematurely when crossed neutral cells that do not change dominate the count; it tests hydrogen only. A starts ionized ($10^{-5}$), which avoids the first & The same test; with a neutral start (\code{xHI\_init = 1}) the outcome depends on how many crossed leaves stay neutral and unchanged, so the share of the first iterations should be checked & Measured, for this case only: clmp1 (\code{xHI\_init = 1}, $10^8$ packets, FUV band on, no-packet share 0.00\%) 99.96\% at iteration 1 (\path{clumpy_test/mochii_clmp1_cell_fraction/run.log}, line 159); clmp2 stopped with He$^0$--He$^+$ measure $d=0.135$\\
11 & An unhit cell keeps its old state: \code{updateCell} copies its gas state and returns (\mc{update_mod.f90}{137--156}) & Every leaf with gas is solved whether or not a packet crossed it (\ms{ion_balance_mod.f90}{1191--1204}; \ms{thermal_mod.f90}{286--306}); the crossing enters only the count (\ms{ion_balance_mod.f90}{1226--1228}; \ms{thermal_mod.f90}{368--370}). An unilluminated leaf can recombine or cool in \MoCHII\ where A keeps it unchanged, which also changes the later field and histories & Not measured; no-packet share 0.00\% in the clumpy runs, so it did not act there\\
\end{longtable}\endgroup

\Measured Before the start of item 9 was matched (helium neutral,
$T_e=10^4$~K, $x_{\rm HI}=10^{-3}$), with the crossing definition above,
the B run clmp2 (\path{clumpy_test/mochii_clmp2_cell_fraction/}, $10^7$
packets in each iteration) stopped at iteration 11 with converged shares
0.43, 0.76, 6.90, 32.28, 45.57, 51.90, 55.74, 59.87, 73.39, 92.86 and
96.14\%. At the stop its He$^0$--He$^+$ change measure was $d=0.135$, and
its [O\,\textsc{iii}]~5007/H$\beta$ was 0.150, against 0.824 after 40
iterations of the same input with \code{conv\_crit='vol'}
(\path{clumpy_test/mochii_clmp2_matched/}). A's run reached 0.0043,
2.07, 71.52, 89.57, 93.01, 94.58 and 95.67\% in its seven iterations
(\path{mocassin/output/summary.out}).

\Source Checked and not acting: A's \code{Xmax = 0.9999} is declared at
\mc{update_mod.f90}{358} and referenced nowhere else in source A, so it
bounds no fraction. \textbf{Not verified:} the size of the effects of
items 2 and 4--7; whether A's first ion-balance call uses the starting
metal fractions as its initial guess (\MoCHII's first solve does not read
them).

\subsection{When the tests answer differently for the same state sequence}
\label{sec:stopcounterexamples}

\Inferred The following are constructed logical examples, not new
MOCASSIN/MoCHII runs. Use the stated defaults unless otherwise specified.
Assume a single grid without echo filtering and that every active cell
receives sampled radiation, so the no-hit exception does not apply. First let 960 of 1000 equal-volume active
cells have no change and a
moving front affect the remaining 40, with $x_{\rm HI}^{\rm old}=0.5$
and $|\Delta x_{\rm HI}|=0.1$ there. A sees 96\% passing cells and meets
its 95\% target even though those front cells fail its 5\% relative test.
MoCHII cell instead sees $\max|\Delta x_{\rm HII}|=0.1>0.001$. If the
new integrated ionized content is $504V$ then
\begin{equation}
 d_{\rm H}=\frac{40V(0.1)}{504V}=7.94\times10^{-3}.
\end{equation}
A similarly persistent front history also exceeds the MoCHII vol
fallback tolerance. If the front occupies much less volume, the vol
answer can differ again. Cell counts, volume integrals, and maxima do
not certify the same spatial accuracy.

\Inferred Conversely, changing every cell from $x_{\rm HI}=10^{-6}$ to
$3\times10^{-6}$ gives A a relative step 2, so its H I percentage test
fails. MoCHII cell sees only $2\times10^{-6}<0.001$ in H II and can pass
if its temperature test also passes. This is a difference in monitored
fraction and normalization, not evidence that one implementation is
numerically malfunctioning.

\Inferred For stationary approximately Gaussian neutral-fraction
increments with standard deviation $\sigma_{\Delta x}$, A's expected
passing-cell fraction is approximately
\begin{equation}
 \Pr(q_i=1)\simeq\operatorname{erf}\!\left(
 \frac{\epsilon_{\rm M}x_{\rm HI}}{\sqrt2\sigma_{\Delta x}}\right).
\end{equation}
It can remain below 95\% with no mean drift. Enabled packet growth may
reduce this noise under the additional finite-variance sampling
assumptions, but its guards and maximum do not guarantee success.
MoCHII's integrated absolute-step numerator likewise has a positive
noise floor; a small signed drift statistic does not cancel that floor
or waive $d<\code{gas\_tol}$. Both can therefore fail their respective
configured tests for a stationary fluctuating state.

\Inferred Slow coherent changes provide the other caution. For a
geometric sequence with $\rho=0.99$ and latest step $5\times10^{-4}$,
the infinite geometric tail would be $0.0495$, far above 0.001. Yet the
code labels the estimate unavailable because $0.99\geq0.95$; its vol
fallback with drift testing off can accept the small latest step after
eight samples. Turning the drift test on rejects such a coherently
signed eight-step history through its large $t_s$. This illustrates the
conditional meaning of the floor branch, rather than a guarantee that
the exact solution error is small whenever an outer flag is true.

\Measured The quoted HII20 C result in Section~\ref{sec:measurements}
uses $2^{25}$ source packets, \code{conv\_crit='vol'}, eight-sample
composition histories, \code{gas\_tol=0.005}, enabled drift testing,
\code{gas\_tol\_te=0.01}, and \code{ion\_relax=1}. Its He$^0$--He$^+$
step floor of about 0.022 exceeds 0.005, so it did not pass that
specified criterion in the 20-pass comparison or before the continuation
was interrupted at 39. The corresponding B floor is about 0.0043.
These are the physical document's Section 5 measurements, not evidence
that a stationary mean is impossible or that another chosen criterion
would necessarily pass. No alternative-criterion run was made here.


\section{MOCASSIN versus v1.10 absorption re-emission}
\label{sec:ac}
\subsection{What becomes closer}
\Source C continues a gas-absorption history immediately at the
absorption point, with a new diffuse energy and isotropic direction, as
ordinary A re-emission does. It retains position/leaf, updates energy,
bin/channel and exact-energy physics, and clears external-origin angular
behavior (\ms{raytrace_amr.f90}{387--400,\allowbreak 479--502};
\ms{diffuse_mod.f90}{575--589}). Part of its diffuse field thus has a
different spatial origin from B's uniform leaf emission.

\Inferred This proximity concerns where and when radiation is emitted.
C remains different: A samples a local combined-emissivity CDF at fixed
history normalization; C uses an absorber-specific yield, changing
weight, and a separately launched explicit remainder.

\subsection{Implemented absorber rule, cap, and explicit remainder}
\label{sec:r7}
\Source The state table has H/He ground continua in channels 1--3,
He I excited/cascade emission in channel 4, and the He III recombination
cascade through He II in channel 5. Channels 6, 7, and optional 8 carry
metal ground continua, free-free, and the selected metal valence-emission
approximation, respectively. These labels specify which emission spectra
are sampled; no atomic-data values are compared here
(\ms{diffuse_mod.f90}{275--332}).

\Source Let $s\in\{\mathrm{H^0,He^0,He^+}\}$ and
$C(s)=\{1\},\{2,4\},\{3,5\}$. In leaf $i$, C freezes
\begin{align}
 N_{{\rm rec},s}&=V\nelec n_{s+}\alpha_{A,s}(\Te),&
 L_s&=\sum_{c\in C(s)}L_c,\\
 N_{{\rm abs},s}&=Vn_s\Gamma_{s,\rm previous},&
 \psi_s&=N_{{\rm abs},s}/N_{{\rm rec},s},\\
 \epsilon_s&=L_s/N_{{\rm rec},s},&
 \epsilon_{{\rm re},s}&=\min(\epsilon_s,E_{{\rm th},s}+\tfrac32k\Te),
 \qquad f_s=\epsilon_{{\rm re},s}/\epsilon_s.
 \label{eq:cap}
\end{align}
The previous raw path-based photoionization rate is multiplied by the
current absorber density. Gas-absorbed energy is not the normalization.
These executable definitions include the cap
(\ms{diffuse_mod.f90}{339--377}; \ms{ion_score_mod.f90}{940--952}). The
balance coefficient defines a recombination count; its numerical data
are outside this comparison.

\Source At incoming photon energy $E$, the absorption continuation yield
and explicit launch are
\begin{equation}
 y_s(E)=\frac{\epsilon_{{\rm re},s}}{E\max(\psi_s,1)},\qquad
 X_c=L_c[1-f_s\min(\psi_s,1)],\quad c\in C(s).
 \label{eq:split}
\end{equation}
The cap removes no state luminosity: its complementary share remains in
$X_c$. Channels 6--8 stay entirely explicit; metal absorption has no
immediate re-emission. Zero recombination count or zero $L_s$ skips
division and retains explicit emission
(\ms{diffuse_mod.f90}{339--377}; \ms{ion_packet_mod.f90}{124--138}).
This is not one fixed-energy packet emitted after each gas absorption.
Enabled He I metastable photoionization has its own rate estimator but
is not a fourth extinction/re-emission absorber in this model; its
coupled state can affect channel 4. The packet gas opacity and the
three $\kappa_s$ refer to the same ground-state absorption terms
(\ms{ion_packet_mod.f90}{92--138};
\ms{ion_score_mod.f90}{485--489}; \ms{diffuse_mod.f90}{303--315}).

\Source C combines scattering and re-emission:
\begin{equation}
 t=\frac{\kappa_{\rm sca}+\sum_s\kappa_s y_s}{\kappa_{\rm ext}},
 \qquad W'=Wt.
 \label{eq:continuation}
\end{equation}
It selects scattering with probability
$\kappa_{\rm sca}/(\kappa_{\rm sca}+\sum_s\kappa_sy_s)$, otherwise absorber
$s$ by $\kappa_sy_s$ and channel by $L_c/L_s$. Re-emitted energy comes
from that channel, not the absorbed energy. $t=0$ terminates; $t>1$ is
allowed without clipping. Roulette follows, and the limit of $10^7$
interactions in one history stops the run through \code{abort\_run}
instead of silently truncating
(\ms{raytrace_amr.f90}{301,\allowbreak 341--421,\allowbreak 505--527};
\ms{diffuse_mod.f90}{575--589}). This is weighted continuation, not a
literal physical photon branching count.

\Source C builds $L_c$, yields, and $X_c$ before clearing the previous
field and before any launch. With no valid field---first pass, restart,
or rebuilt grid---all yields stay zero and full state emission is
separately launched. Setup/reset invalidates the field; reduction
publishes the next valid field (\ms{ionizing_field_mod.f90}{119--132};
\ms{diffuse_mod.f90}{340--348};
\ms{ion_score_mod.f90}{272--338,\allowbreak 345--411,\allowbreak 831--833,\allowbreak 940--952}). Explicit
remainder packets themselves follow the same gas rule and can produce
further descendants. A re-emitted packet stays in the interaction loop;
its next flight is sampled, rather than starting another analytic
launch-ray integration (\ms{raytrace_amr.f90}{395--400,\allowbreak 464--466}). Dust-IR packets use their distinct band
(\ms{raytrace_amr.f90}{315}; \ms{ionizing_field_mod.f90}{175--186}).

\subsection{Remaining differences: conservation and estimators}
\label{sec:conservation}
\Inferred For an active absorber with a valid previous field, positive
$N_{{\rm rec},s}$ and positive $L_s$, freeze coefficients and define
$b_c=f_sL_c/\max(N_{{\rm rec},s},N_{{\rm abs},s})$. The expected
mean re-emitted power is $b_c\mathbb E[\widehat N_{{\rm now},s}]$,
where $\widehat N_{{\rm now},s}$ is the current pass absorption-count
estimator. For completed histories with finite first moments,
\begin{align}
 \mathbb E[\widehat R_c-X_c-b_c\widehat N_{{\rm now},s}]&=0,
 \label{eq:countidentity}\\
 X_c+b_cN_{{\rm abs},s}&=L_c.
 \label{eq:fixedidentity}
\end{align}
Here $\widehat R_c$ is total explicit plus re-emitted channel power.
Inactive absorbers have $b_c=0$ and $X_c=L_c$, without division.
The second identity is algebraic; equating current absorption counts to
frozen previous-field counts requires a normalization fixed point. It
does not prove coupled gas convergence or imply
$\mathbb E[\widehat R_c]=L_c$ in arbitrary transient passes. An unrelaxed photoionization-equilibrium solve
(\code{ion\_relax=1}), with consistent recombination counts and no
other ionization, gives $\psi=1$; secondary ionization, charge exchange,
and metastable coupling prevent a universal $\psi=1$ claim.

\Source C preserves the existing gross band-output meaning:
\begin{align}
 L_{\rm EMIT}&=L_{\rm source}+L_{\rm LEAFEM}+L_{\rm REEMIT}
              =L_{\rm ABS}+L_{\rm ESC}+L_{\rm LOST},\\
 L_{\rm NETABS}&=L_{\rm ABS}-L_{\rm REEMIT}.
\end{align}
\code{L\_LEAFEM} is explicit launched power, not all state emissivity.
Deposition $W-W'$ is signed at $t>1$; realized roulette also has a signed
correction. The ledgers retain these terms and count photon rates as
luminosity divided by current energy
(\ms{ion_score_mod.f90}{637--737,\allowbreak 1090--1116,\allowbreak 1152--1182}).
\code{GasAbsorbedPower} is a path estimate of gross gas absorption, not
the r7 yield input (\ms{ion_score_mod.f90}{492--498,\allowbreak 940--952}).

\Inferred A keeps energy normalization in ordinary re-emission while
changing represented photon number. C keeps weighted accounting and
matches state-emission normalization at its count fixed point. B injects
state-derived diffuse power independently of individual absorptions.
None conserves physical photon number across ionization/recombination
and redistribution. Computational band closure also does not prove
bolometric thermal balance, especially before gas/dust convergence or
outside the transported band.

\Source C retains B's analytic direct flight, continuous energies,
exact-energy gas estimators, scheduling, QMC launch option, and gas
stopping criteria. Its observer pass adds isotropic re-emission peel
with continued weight before roulette, including dust-free inputs and
formerly external photons (\ms{raytrace_amr.f90}{395--400};
\ms{ion_peel_mod.f90}{276--285};
\ms{ionizing_field_mod.f90}{121--126}). It does not replace these with
A's frequency-bin estimator or escaped-packet binning.

\section{The v1.00 and v1.10 choices directly compared}
\label{sec:bc}
\subsection{Spatial model, diffuse sample size, and warm starts}
\Inferred B distributes leaf state emission uniformly in its volume;
C places the recycled part where absorptions occur and leaves the
remainder uniform. Radiation and absorption density need not be uniform
inside a cell with constant gas coefficients. These are distinct
finite-grid spatial models, not merely different samplers of one
cell-volume prescription. Refinement can reduce their difference;
coarse-grid equality is not required by the count identity. The position
choices are \ms{diffuse_mod.f90}{504--507,\allowbreak 686--704} and
\ms{raytrace_amr.f90}{479--502}.

\Inferred In A the ordinary diffuse-event count follows absorption
histories and terminal/escape outcomes, with extra diffuse/resonance
populations when enabled. In B the independent diffuse count follows
full state luminosity in Equation~\eqref{eq:ndif}. In C that formula
uses $\sum X_c$ and other diffuse samples are absorption descendants.
Fewer explicit packets can reduce independent diffuse sampling and
increase correlation among descendants of the same launch. Immediate
re-emission therefore does not guarantee less noise or faster execution
at equal source counts. Initial explicit launches can use QMC, but a
variable-length re-emission chain is not thereby fixed-dimensional QMC.

\Source A new-mode warm start restores physical state but resets
stopping histories and begins with full explicit emission. Refinement
invalidates the estimator epoch and clears stopping histories
(\ms{previous_state_mod.f90}{296--307,\allowbreak 349--377};
\ms{main.f90}{103--114,\allowbreak 264--275}; \ms{ion_score_mod.f90}{272--343}).
B can retain optional histories. No old energy-normalization state is
needed or restored for the r7 rule.

\subsection{One packet history in parallel}
\label{sec:histories}
\Source Table~\ref{tab:histories} follows the connected launch,
transport, and emission calls cited in Sections~\ref{sec:flight},
\ref{sec:diffuseb}, and \ref{sec:r7}. Ordinary line termination and the
separate dust-IR calculation must not be confused with continuing
gas-band histories.
\begin{table}[p]
\centering\small\setlength{\tabcolsep}{5pt}
\begin{tabular}{@{}P{.305\linewidth}P{.305\linewidth}P{.305\linewidth}@{}}
\toprule A: MOCASSIN & B: leaf emissivity & C: absorption re-emission\\\midrule
\begin{minipage}[t]{\linewidth}\vspace{0pt}
\begin{enumerate}[leftmargin=1.35em,itemsep=6pt]
\item Take the source normalization, choose a discrete frequency, and draw direction.
\item An ordinary line terminates. Continuum below the H edge escapes immediately without dust.
\item Draw exponential optical depth; walk cell faces and score frequency-indexed path lengths.
\item At escape, add energy to frequency/angle/origin bins and end.
\item If dust scattering is chosen, preserve normalization/frequency, change direction, and repeat step 3.
\item At absorption, keep the actual position and normalization. Select the cell's combined diffuse or dust-only spectrum.
\item Initialize the new frequency/direction/type and return to step 2. Solve gas/dust after the population.
\end{enumerate}\end{minipage}
&\begin{minipage}[t]{\linewidth}\vspace{0pt}
\begin{enumerate}[leftmargin=1.35em,itemsep=6pt]
\item Launch continuous-energy sources and separate full-state diffuse packets. Each starts with $w=1$.
\item Add direct peel if requested, then score the analytic attenuated launch ray and direct escape.
\item Without dust scattering, end; absorption already contributed to path/rate estimates.
\item Force a first interaction with extinguished weight and conditional depth; do not count its incoming path twice.
\item Multiply by the dust-scattering albedo, peel, and apply low-weight roulette.
\item Choose HG direction and exponential next depth; score actual traversed segments.
\item Repeat until escape/termination. No immediate gas-recombination descendant is made.
\end{enumerate}\end{minipage}
&\begin{minipage}[t]{\linewidth}\vspace{0pt}
\begin{enumerate}[leftmargin=1.35em,itemsep=6pt]
\item Freeze $L_c,y_s,X_c$ before launch. Invalid old field means $y=0$, $X=L$.
\item Launch sources/explicit remainder and score the same analytic direct flight as B.
\item Force an interaction at positive extinction, also without dust; exclude duplicate incoming path scoring.
\item Compute $t$. End at zero; otherwise multiply weight by $t$, without clipping.
\item Select dust scattering or absorber by $\kappa_sy_s$. Re-emission keeps position, samples channel/energy and isotropic direction, rebuilds physics, and peels before roulette.
\item Add dust-scattering peel when selected, before roulette; choose HG direction for scattering. Draw an exponential flight and score segments.
\item Repeat until termination. Explicit descendants follow the same rule; solve gas after the pass.
\end{enumerate}\end{minipage}\\\bottomrule
\end{tabular}
\caption{Side-by-side pseudocode. A keeps normalization within an
ordinary history; B/C analytically integrate the first flight and use
weighted continuation. This is source-read control flow, not timing evidence.}
\label{tab:histories}
\end{table}

\FloatBarrier
\subsection{Existing measurements: only B versus C}
\label{sec:measurements}
\Measured Values here are quoted from \path{docs/MoCHII_physics.tex},
Section 5 (G-R6, Lexington, convergence/noise, and P3), and
\path{md/ABSORPTION_REEMISSION_GATES_REPORT_claudeB.md}, sections
``r7 with the writer,'' ``Round 3,'' ``Gates on the final binary,'' and
subsequent documentation results.
No data were regenerated and no MOCASSIN run was compared. The table
states the input scope; coefficient and line-list differences are excluded.

\begin{table}[htbp]\centering\small\setlength{\tabcolsep}{4pt}
\begin{tabular}{@{}P{.18\linewidth}P{.38\linewidth}P{.38\linewidth}@{}}
\toprule Condition & G-R6 sphere & Lexington HII40 / HII20\\\midrule
Geometry & 8 pc box, gas sphere radius 4 pc; uniform levels 4--6 & Matched level-6 grids in each mode; 262144 leaves\\
Gas & Pure H, 100 cm$^{-3}$; fixed $10^4$ K & Same H/He/metals and density grid in each mode; temperature solved within 3000--30000 K\\
Source & 40000 K Planck spectrum; $Q=10^{49}$ s$^{-1}$ & 40000 / 20000 K; matched input ionizing powers $1.2734\times10^{39}$ / $2.5326\times10^{38}$ erg s$^{-1}$\\
Radiation & Matched transported continuum/Milne sampling; local Lyman-line treatment & 13.598--100 eV with same FUV extension; 32 ionizing bins\\
Sampling & $2^{20}$ Sobol source packets; $K=2,4,2$ at levels 4,5,6 & $2^{25}$ source packets; identical seed 20260728\\
Iteration & At most 80 passes; stated sphere criterion & Same m13 state; volume/drift criterion; composition tolerance 0.005 and thermal tolerance 0.01\\
Caveat & Some warm-started realizations; level-6 error is a lower bound & HII40 both pass the specified test; HII20 C does not pass it in 20 fixed iterations\\\bottomrule
\end{tabular}
\caption{Matched conditions quoted from physical-document Section 5.
No MOCASSIN column is implied. These are static transfer problems
without gas dynamics or gravitational evolution.}
\end{table}

\FloatBarrier
\paragraph{The stopping criterion used for the quoted Lexington runs.}
\Source The actual saved \path{run.in} files in
\path{scratch_v1.10_claude/lexington/} specify
\code{conv\_crit='vol'}, \code{conv\_drift\_test=.true.},
\code{gas\_tol=0.005}, \code{gas\_tol\_te=0.01},
\code{ion\_relax=1}, and 33554432 source packets
(the HII20 C 20-pass input, lines 7 and 31--36). Both modes use the
same physical-state restart. The following is the implemented test at
these settings, rather than a claim of universal convergence:
\begin{enumerate}[leftmargin=1.7em,itemsep=2pt]
\item Require eight samples for each active composition history:
H$^+$, He$^0$--He$^+$, and He$^+$--He$^{2+}$ only when integrated
He$^{2+}$ exceeds 1\% of integrated ionized helium.
\item If the fitted step-decay ratio is below 0.95 and its distance
estimate is available, require remaining distance below 0.005.
Otherwise require the latest normalized composition step below 0.005
\emph{and} the absolute signed-mean step divided by its standard error
below 2.365 over the eight-sample window.
\item Also require
$\sum V n_{e,\rm new}|\Delta T_e|/\sum V n_{e,\rm new}T_{e,\rm new}<0.01$.
\end{enumerate}
The H and He$^0$ steps are
$\sum V|\Delta x_{\rm HII}|/\sum Vx_{\rm HII,new}$ and
$\sum V|\Delta x_{\rm HeI}|/\sum V(x_{\rm HeII}+x_{\rm HeIII})_{\rm new}$,
respectively. These are integrated measures, not maximum local changes
(\ms{main.f90}{503--508,\allowbreak 545--550};
\ms{ion_balance_mod.f90}{1225--1254,\allowbreak 1340--1349};
\ms{thermal_mod.f90}{366--373,\allowbreak 426--431};
\ms{fixed_point_distance_mod.f90}{69--74,\allowbreak 119--192}).
The C input for the tabulated HII20 comparison caps iteration at 20
and permits output without passing this test (saved input, lines 31
and 46). The separate continuation retained the criterion and was
interrupted at iteration 39; it did not complete its nominal
100-iteration cap. This is a result for these sampling and stopping
settings, not evidence that the method cannot approach a stationary
mean under other criteria.

\Measured Sphere-radius differences $(R_C-R_B)/R_B$ are
$-0.059\%\pm0.003\%$, $-0.039\%\pm0.002\%$, and
$+0.001\%\pm0.006\%$ at levels 4, 5, 6. These are quoted errors, not
new intervals. Some samples continued converged states; the reported
level-6 error is a lower bound. At level 5 the emission-normalization
residual is $+4.9\times10^{-4}\pm3.7\times10^{-5}$, resolved but within
the stated $5\times10^{-3}$ allowance. Frozen-state G-R4 reports a thick-case
absorbed-power difference falling from about 2.3\% at level 0 to 0.18\%
at level 3 with $10^5$ packets, with similar trend at $4\times10^5$.
These values are from the binary of 11:07. On the final binary of
2026-11-13 (20:55) the gates were rerun (physical document, Section~5.4, paragraph ``The gates on the final
binary''): the
sphere differences are $-0.062\%$, $-0.037\%$, $-0.002\%$ at levels 4, 5,
6 (16 converged spheres continued by restart), the level-5
normalization residual $+4.68\times10^{-4}\pm1.2\times10^{-5}$, and the
G-R4 thick difference 2.12\% at level 0 and 0.17\% at level 3; one G-R7
row (a channel carrying $2\times10^{-4}$ of the emission) missed its
stated resolution limit at four realizations without a detected bias.
These concern finite-grid spatial models, not accuracy against A.

\Measured Lexington $\langle T[N_pN_e]\rangle$ means are 8176.13/8169.62 K (B/C, HII40)
and 6942.21/6938.00 K (HII20); H$\beta$ luminosity ratios C/B are
0.99960 and 0.99996. Selected neutral-edge line outputs change by about
+3.9\%/+6.1\% for [O I] 6302 and +2.2\%/+4.3\% for [N I] 5200
(HII40/HII20). They illustrate the algorithm comparison, not atomic-data
or line-list differences. With the criterion specified above, HII20 C did not pass in the
20-iteration comparison or before its continuation was interrupted at
iteration 39. Its quoted helium step near 0.022 exceeds the 0.005
composition threshold in the noise-floor branch, while the temperature
step is below the separate 0.01 threshold. It therefore lacks a stopping
certificate under those conditions; this is not a claim of unbounded
drift or nonexistence of a stationary mean.

\Measured Final explicit diffuse counts fall from 19188289 to 69064
(HII40) and 16935338 to 18954 (HII20). Volume He$^0$--He$^+$ step floors
rise from $1.3\times10^{-3}$ to $4.3\times10^{-3}$ (HII40), and
$4.3\times10^{-3}$ to $2.1$--$2.2\times10^{-2}$ (HII20). These are
iteration fluctuations under stated inputs, not independent measurements
of universal variance. The reports attribute the higher stationary noise
to fewer explicit diffuse packets and sampled re-emission chains; they
did not isolate contributions of individual random choices. No
noise-reduction change is assumed here.

\Measured P3 quotes eight realizations at two source counts in a frozen
helium-containing shell. Four global variance ratios fit the broad
stated $1/N$ test, but the original all-leaf He I rate-variance criterion
fails in both modes. A later resolved-leaf selection gives a different,
explicitly post-selected result. Maximum observed history weights and
those selected samples do not establish a general $t>1$ variance bound
(physical document, Section 5, ``Variance (P3)''; gate report, ``Round 3'').

\Measured Existing G-R10 imaging passes under the separately predeclared
transparent-source quadrature subtraction. Without that subtraction the
round-4 external-dust 32/128-direction bound is 3.21\%, above the 3\%
margin, and the round-3 raw FAIL stands as the record. G-R11 includes 16
reconverged pure-H restart pairs; maximum complete-interval state/rate
bounds are 0.218488\%/0.227666\%, inside 1\%/3\%. These are quoted from
\path{md/ABSORPTION_REEMISSION_ROUND4_codex.md} and its protocols,
summarized in physical-document Section 5. They do not cover arbitrary
helium/metal weight-gain histories.

\Inferred The spatial trend supports the finite-grid interpretation on
those inputs. The quoted noise/counts show why immediate re-emission
alone promises no noise improvement. Neither proves universal
convergence, a continuum-limit theorem, or a ranking against MOCASSIN.

\begin{figure}[p]\centering
\includegraphics[width=.86\textwidth,height=.77\textheight,keepaspectratio]{figures/reem_g1_profiles.pdf}
\caption{Existing G-R6 full radial profiles at levels 4--6, reproduced
without recalculation from physical-document Section 5. Dashed: B;
solid: C. H I/electron-density profiles and differences are shown.
This is a quoted B/C result, not an A/B/C calculation. Conditions and
statistical limitations are in Section~\ref{sec:measurements}.}
\label{fig:g1}\end{figure}
\begin{figure}[p]\centering
\includegraphics[width=.82\textwidth,height=.77\textheight,keepaspectratio]{figures/reem_hii40_profiles.pdf}
\caption{Existing Lexington HII40 full profiles from physical-document
Section 5: H/He composition, temperature, electron density, and B/C
differences. Both runs passed the volume/eight-step/drift criterion at composition
tolerance 0.005 and thermal tolerance 0.01, with 33554432 source packets. No MOCASSIN comparison
or new calculation is represented.}\label{fig:hii40}\end{figure}
\begin{figure}[p]\centering
\includegraphics[width=.82\textwidth,height=.77\textheight,keepaspectratio]{figures/reem_hii20_profiles.pdf}
\caption{Existing Lexington HII20 full profiles from physical-document
Section 5. C is the 20-iteration result. With 33554432 source packets, volume
stopping, an eight-step composition window at tolerance 0.005, the
drift test, and thermal tolerance 0.01, C does not pass that specified
criterion; see Section~\ref{sec:measurements}. Integrated temperature or H$\beta$ agreement alone does not
establish profile agreement. No new calculation is represented.}
\label{fig:hii20}\end{figure}
\begin{figure}[p]\centering
\includegraphics[width=.96\textwidth,height=.68\textheight,keepaspectratio]{figures/reem_iteration_history.pdf}
\caption{Existing physical-document Section 5 iteration histories,
reproduced without recalculation. Composition and thermal steps have
different fluctuation floors. With 33554432 source packets and the
volume/eight-step/drift criterion, HII20 C has a He step near 0.022,
above composition tolerance 0.005 (thermal tolerance 0.01), so it does
not pass that specified test.
These finite runs do not isolate individual noise sources.}
\label{fig:noise}\end{figure}

\clearpage
\section{What was not verified}
\label{sec:unknown}
The following eight items delimit this comparison.
\begin{enumerate}[leftmargin=2em,itemsep=5pt]
\item \textbf{Independent release authentication.} Source A is the
original-author tree identified by the author on November 14, 2026,
including the five supplied \path{~/MOCASSIN/error_fix/} files verified
identical to its source copies. The separate \path{~/RT_Codes/} tree is
the user's unfinished correction/optimization work. Neither an official
release archive nor the authorship of every dated correction was
independently authenticated.
\item \textbf{Matched numerical cross-code performance.} No MOCASSIN
binary was built or run. Accuracy, time, memory, convergence rate, energy
budgets, and efficiency of A versus B/C were not numerically compared
with identical inputs.
\item \textbf{General C convergence/variance.} No proof of coupled
gas/transport convergence or finite variance for arbitrary $t>1$
histories follows from r7 identities/gates. Cross-absorber returns and
hot or strongly coupled states remain outside such a claim.
\item \textbf{Converged HII20 C or a noise remedy.} The quoted HII20 C
run does not pass the volume/eight-step/drift test at composition tolerance
0.005 and thermal tolerance 0.01 with 33554432 source packets, over
20 tabulated iterations or its continuation interrupted at 39.
Additional counts, different stopping,
and noise-reduction methods were not implemented/tested here.
\item \textbf{Broader restart/imaging coverage.} Quoted pure-H restart
and H/dust imaging do not certify full helium/metals, large gains, or
every geometry/parallel combination.
\item \textbf{Residual offsets and Lexington grid limits.} The small
count-residual offsets of G-R9 were positive on the 11:07 binary and of
both signs on the final binary; four realizations a case do not establish
them as bias or noise. Shrinkage
of Lexington neutral-edge differences under refinement is unmeasured.
\item \textbf{All optional MOCASSIN outputs/warm routes.} Connected
ordinary transfer, resonance approximations, output, and warm paths were
read, but their numerical correctness, zero-depth behavior, exact
random-state continuation, and rank-count reproducibility were not
executed or certified.
\item \textbf{The binned source spectrum of MOCASSIN.} The quoted
bin counts and photon-rate ratios were previously reconstructed from the
input spectrum and a frequency grid in the student's clumpy-model
output; no such reconstruction or transport was rerun here. According to
the author, that run used the original-author code. The accessible
student tree \path{/home/dg.a9mon/test_Mocassin/mocassin_fix2/source/}
matches A in \code{mocassinPlot}, \code{grid\_mod}, \code{output\_mod},
\code{continuum\_mod}, \code{photon\_mod}, and \code{emission\_mod};
\code{mocassinPlot} writes the same four-integer plot-grid header (\mc{mocassinPlot.f90}{474--482}). Its
\code{iteration\_mod} includes the April 16, 2025 column-wise MPI
changes, \code{update\_mod} uses a thermal cap 25 rather than A's 50 and
initializes photoionization sums differently, and \code{set\_input\_mod}
uses a different output-directory inquiry. Thus the student tree is not
identical to A as a whole. The \code{rauch} binning expressions are the
same (\mc{continuum_mod.f90}{216--283,\allowbreak 442--458}); unlike
\code{SB99}, they do not interpolate empty bins. The executable used
for the existing output was not forensically identified, and its binned
spectrum was not printed by a new run. The share of the model's
metal-line differences caused by this source mismatch was not computed.
\end{enumerate}

\paragraph{Source caveats found during inspection.}
\Source The following original-source expressions deserve separate notice.
\code{writeContCube} computes a requested frequency range but sums
\code{freq=1,nbins}, ignoring the selection
(\mc{output_mod.f90}{2756--2763,\allowbreak 2777--2780}). The resonance escape formula
allows $\tau_\mu=0$, then evaluates
$[1-\exp(-2\tau_\mu)]/(2\tau_\mu)$: its continuous limit is 1 but its
literal zero argument is $0/0$ (\mc{emission_mod.f90}{3051--3065}).
These are source findings, not runtime measurements; reachability and
consequences of the second were not tested.
The qHeat routine additionally zeros its probability array, includes
all temperature states in the normalization, but fills only the first
\code{nTbins-1} probabilities. Thus the returned sum lacks the final
state's share; an unsupported-grain return can also leave the
temperature-grid and probability arrays zero (\mc{emission_mod.f90}{1465--1488,\allowbreak 1660--1677}). These are
source-level limitations of the optional path, not measured effect
sizes or new convergence results. Its ordinary outer percentage does
not certify stochastic-temperature-distribution accuracy.
 External sources were
preserved. Neither output is a validated counterpart of MoCHII peel or
a performed MOCASSIN comparison.

\end{document}
